\section{引言}\label{sec:intro}

场论中的梯度流已被证明是一个非常有用的概念。它最初作为正则化格点计算中
发散与紫外涨落的方法被提出\,\cite{Narayanan:2006rf,Luscher:2009eq}，
随后主要凭借其在格点 \qcd{} 标度设定（scale-setting）中的应用而获得广泛
关注\,\cite{Luscher:2010iy,Luscher:2013vga,Borsanyi:2012zs}，如今已成为该
领域许多实际计算中不可或缺的工具（例如文献
\citeres{Luscher:2013cpa,Sommer:2014mea,Ramos:2015dla,Aoki:2019cca}）。
尽管最初聚焦于 \qcd{}，梯度流形式（\gff）的推广带来了广泛得多的应用范围，
例如场论中对偶性的研究\,\cite{Aoki:2015dla,Aoki:2016ohw,Aoki:2017uce,Aoki:2017bru,
  Aoki:2018dmc,Aoki:2019bfb}。

在 \qcd{} 中，将其表述为五维场论的方案见于 \citere{Luscher:2011bx}，其中
额外维度与所谓的流时间相联系。基本理论（即真实的 \qcd{}）的形式充当流时间
为零处的边界条件。已经证明：当用重正化后的参数与场表达时，正流时间处的复合
算符是有限的\,\cite{Luscher:2011bx}。这一性质，加上认识到梯度流可以用作格点
计算与微扰计算之间的匹配方案，照亮了一整套正在被探索的全新应用。

微扰 \qcd{} 与格点 \qcd{} 之间最早可能的相互滋养之一，是利用有限流时间处的
胶子凝聚（\qcd{} 作用量密度）定义强耦合的新方案\,\cite{Luscher:2010iy}。
它与 \msbar{} 耦合的微扰关系已知至三圈\,\cite{Harlander:2016vzb}。
\footnote{除非另有说明，本文中的微扰计算均指无限体积下的计算。纳入有限体积
效应需要不同的技术，例如数值随机微扰理论（Numerical Stochastic Perturbation
Theory），见 \citere{DallaBrida:2017tru}。}另一个跨越边界（cross-boundary）
的应用是能量动量张量的梯度流定义，它利用小流时间展开，把能量动量张量表示为
正流时间处良定的复合算符与微扰可及的系数函数\,\cite{Suzuki:2013gza,Makino:2014taa,
  Harlander:2018zpi,Iritani:2018idk}。再一个例子是通过使用梯度流把欧氏准
\pdf{}（quasi PDF）\footnote{\pdf=parton distribution function，即部分子分布
函数。}与微扰光锥 \pdf{} 联系起来的提议\,\cite{Monahan:2016bvm,Monahan:2017hpu}。

所有这些应用都同时依赖来自格点与来自微扰计算的输入，且二者都在有限流时间处。
人们发现，微扰论中的高阶修正对于减小以重正化标度变化估计的微扰截断误差至关
重要\,\cite{Harlander:2016vzb,Harlander:2018zpi}。这类修正还被发现能显著稳定
相应格点结果所需的外推\,\cite{Iritani:2018idk}。

尽管 \gff\ 中的圈积分包含额外的指数因子以及对流时间变量的积分，正则微扰计算
的许多重要技术仍保持其效用，只是形式略有推广。本文的目标是为 \gff\ 内的进一
步微扰计算提供基础，从而为进一步探索这一方法的能力作出贡献。\sct{sec:pert}
回顾微扰论中的 \qcd{} 梯度流，概述 \citeres{Luscher:2010iy,Luscher:2011bx,
Luscher:2013cpa} 的结果。关联函数多圈微扰计算自动化的各个阶段在
\sct{sec:implementation} 中描述：费曼图的生成、费曼规则的插入与求值，随后把
圈积分约化到一组主积分，最后对后者进行数值计算。作为应用，\sct{sec:observables}
给出胶子凝聚的三圈结果（先前已在 \citere{Harlander:2016vzb} 中获得）、夸克凝
聚的结果，以及属于新结果的``加环''夸克场重正化常数。这些量使我们能够定义精
确的梯度流耦合与梯度流质量方案及其到 \msbar{} 的匹配，如 \sct{sec:couplingandmass}
所示。我们的结论见 \sct{sec:conclusions}。
